\id{ҒТАМР 65.63.33}{}

\begin{header}
\swa{А.Ж. Оразов, Б.Ж. Рыскалиева, А.А. Сармалаев, А.М. Меркенова}{KF ҰЙЫТҚЫ ДАҚЫЛЫМЕН ӘРТҮРЛІ СҮТ ТҮРЛЕРІНІҢ АШУ ПРОЦЕСІН ЗЕРТТЕУ}

А.Ж.~Оразов\envelope,
Б.Ж.~Рыскалиева,
А.А.~Сармалаев,
А.М.~Меркенова
\end{header}

\begin{affil}
Жәңгір хан атындағы Батыс Қазақстан аграрлық-техникалық университеті, Орал, Қазақстан

\corrauthor{Корреспондент-автор: orazov\_ayan@mail.ru}
\end{affil}

Мақалада сиыр, ешкі және түйе сүттерін KF ұйытқы дақылымен ашыту
процесіне салыстырмалы талдау жасалды. Зерттеу барысында әртүрлі сүт
шикізатынан дайындалған айран үлгілерінің органолептикалық және
физика-химиялық көрсеткіштері, сондай-ақ титрленетін және белсенді
қышқылдық динамикасы зерттелді. Алынған нәтижелер бойынша KF ұйытқысы
барлық сүт түрлерінде сүт қышқылы бактериялары мен ашытқылардың белсенді
дамуын қамтамасыз етіп, өнімнің 90-120°Т қышқылдық деңгейіне жетуіне
мүмкіндік берді. Алайда, ашыту процесінің ұзақтығы шикізат түріне
тікелей байланысты екені анықталды: сиыр және ешкі сүті үшін ашу уақыты
4-5 сағатты құраса, түйе сүтінде бұл процесс 8 сағатқа дейін созылды.
Бұл кешеуілдеу түйе сүтінің құрамындағы ақуыздар мен майлардың жоғары
мөлшерімен және оның биохимиялық ерекшеліктерімен түсіндіріледі. Зерттеу
қорытындысы бойынша KF ұйытқы дақылы функционалдық қасиеттері бар сүт
өнімдерін өндіруде тиімді деп табылды, әрі түйе сүтінен өнім алу үшін
технологиялық режимді бейімдеу қажеттілігі негізделді.

{\bfseries Түйін сөздер:} ашыту, сүт, айран, ұйытқы, қышқылдық, құрылым.

\begin{header}
ИССЛЕДОВАНИЕ ПРОЦЕССА ФЕРМЕНТАЦИИ РАЗЛИЧНЫХ ВИДОВ МОЛОКА ЗАКВАСОЧНОЙ КУЛЬТУРОЙ KF

А.Ж.~Оразов\envelope,
Б.Ж.~Рыскалиева,
А.А.~Сармалаев,
А.М.~Меркенова
\end{header}

\begin{affil}
Западно-Казахстанский аграрно-технический университет им. Жангир хана, Уральск, Казахстан,

е-mail: orazov\_ayan@mail.ru
\end{affil}

В статье представлены результаты сравнительного анализа процесса
ферментации коровьего, козьего и верблюжьего молока с использованием
заквасочной культуры KF. Целью работы являлась оценка возможности
использования различных видов молочного сырья для производства
кисломолочного продукта (айрана) и изучение качественных характеристик
полученных образцов. В ходе исследования были изучены органолептические
и физико-химические показатели, а также динамика титруемой и активной
кислотности. Установлено, что применение культуры KF обеспечивает
активное развитие молочнокислых бактерий и дрожжей во всех образцах,
способствуя формированию устойчивой консистенции, специфического вкуса и
аромата. Кислотность готовых продуктов варьировала в пределах 90-120 °Т.
Выявлено, что продолжительность процесса сквашивания существенно зависит
от вида сырья: ферментация коровьего и козьего молока завершалась за 4-5
часов, тогда как сквашивание верблюжьего молока длилось до 8 часов.
Данная особенность обусловлена высоким содержанием белков и жиров, а
также специфическим биохимическим составом верблюжьего молока.
Результаты исследования подтверждают эффективность использования
закваски KF для разработки технологии функциональных кисломолочных
продуктов, однако указывают на необходимость адаптации временных режимов
для верблюжьего молока.

{\bfseries Ключевые слова:} ферментация, молоко, айран, закваска,
кислотность, структура.

\begin{header}
INVESTIGATION OF THE FERMENTATION PROCESS OF VARIOUS TYPES OF MILK WITH STARTER CULTURE KF

A.Zh.~Orazov\envelope,
B.Zh.~Ryskaliyeva,
A.A.~Sarmalaev,
A.M.~Merkenova
\end{header}

\begin{affil}
Zhangir khan West Kazakhstan Agrarian and Technical University, Uralsk, Kazakhstan,

е-mail: orazov\_ayan@mail.ru
\end{affil}

This article investigates the influence of the fermentation process on
various types of milk (cow, goat, and camel) using the KF starter
culture. The study evaluates the feasibility of using these milk types
to produce a fermented milk product (ayran) and analyzes the
organoleptic and physico-chemical parameters of the resulting samples.
By observing the dynamics of titratable and active acidity, the optimal
fermenta\-tion regime for this starter culture was determined. The results
indicate that the acidity of the products made from cow, goat, and camel
milk ranged from 90 to 120 °T. It was observed that the product made
from camel milk possessed lower acidity and differed significantly in
taste and texture. The application of the KF starter culture ensured the
active development of lactic acid bacteria and yeast, contributing to
the formation of a stable consistency, distinct taste, and aroma. The
duration of the fermentation process varied depending on the raw
material: while cow and goat milk required 4--5 hours, the fermentation
of camel milk lasted up to 8 hours. This delay is attributed to the
relatively high protein and fat content and the specific biochemical
composition of camel milk. The study concludes that the KF starter
culture is effective for the production of fermented milk products with
improved functional properties, though specific adjustments to the
fermentation time are required for camel milk.

{\bfseries Keywords:} fermentation, milk, ayran, starter culture, acidity,
structure.

\begin{multicols}{2}
{\bfseries Кіріспе.} Ашымалы сүт өнімдерінің диеталық қасиеттері асқазан
мен ішек жолдары белсенділігіне әсер етеді, нәтижесінде ас қорыту
жолдары тағамның қорытылуын тездететін ферменттерді қарқынды шығарады.
Жағымды, сәл сергітетін және өткір дәмі бар ашымалы сүт өнімдері тәбетті
ашады, ас қорыту жолдарының белсенділігін арттырады, осылайша ағзаның
жалпы жағдайын жақсартады. Сүт өнімдерінде кездесетін микроорганизмдер
тобы, ферментациялық процестерге қатысып, лактозаны сүт қышқылына
айналдырады. Бұл бактериялар, негізінен, \emph{Lactobacillus,
Bifidobacterium,} \emph{Streptococcus} және \emph{Lactococcus} тектес
микроорганизмдер қатарына кіреді және олардың адам ағзасына пайдалы
әсерлері бар {[}1{]}. Сүт қышқылды бактериялардың (СҚБ) маңызды
қасиеттерінің бірі --- олар ас қорыту жүйесін сақтай отырып,
микрофлораны қалыпқа келтіреді, яғни пайдалы бактериялардың көбеюіне
ықпал етеді және зиянды микроағзалардың өсуін тежейді. Сонымен қатар,
СҚБ-лар иммундық жүйенің күштілігін арттыруға, ішек микрофлорасының
балансын сақтауға, ас қорытуды жақсартуға және тағамның қоректік
құндылығын арттыруға көмектеседі {[}2{]}.

Сиыр сүті ақуыздардың, майлардың және көмірсулардың теңдестірілген
құрамына байланысты сүт өнеркәсібінде ең көп таралған шикізат түрі болып
табылады. Оның құрамында барлық қажетті қоректік заттар бар және ол адам
ағзасы үшін өте маңызды тағамдардың бірі. Сүттің ақуыздары негізінен
казеин мен сарысу ақуыздарынан тұрады, олар ағзаның өсуі мен қалпына
келуі үшін қажет аминқышқылдарына бай {[}3{]}.

Ешкі сүті өзінің жоғары сіңімділік қасиеттерінің арқасында сиыр сүтіне
қарағанда адам ағзасы үшін анағұрлым қолайлы болып табылады, бұл оны ас
қорыту жүйесі бұзылған адамдар үшін тиімді тағамдық ресурс ретінде
қарастыруға мүмкіндік береді {[}4,5{]}. Статистикалық деректерге сәйкес,
ешкі сүтінің әлемдік тұтыну көлемі сиыр сүтінен кейінгі екінші орынды
иеленеді {[}6{]}. Құрамы мен функционалдық қасиеттері бойынша ешкі
сүтінен алынатын өнімдер сиыр сүтінен өндірілетін өнімдерге балама
ретінде жоғары технологиялық және тағамдық маңыздылыққа ие. Осыған
байланысты ешкі сүтін қолдана отырып, емдік-профилактикалық бағыттағы
ашытылған сүт өнімдерінің технологиясын әзірлеу қазіргі заманғы сүт
өнімдерін өндіру саласындағы өзекті ғылыми-зерттеу және практикалық
міндеттердің бірі болып саналады {[}7,8{]}.

Елімізде сиыр және ешкі сүттерімен қатар балама ретінде түйе сүтінен
алынатын өнімдер кеңінен тұтынылады. Бұл өнімдерге сұраныс әсіресе
Орталық Азия, Африка және Таяу Шығыс елдерінің тұрғындары арасында
климаттық ерекшеліктеріне байланысты жоғары. Ол тағамдық және
биологиялық құндылығы тұрғысынан, шөлейт және аридті климаттық
белдеулерде өмір сүретін халықтар үшін маңызды ақуыз көзі болып
табылады. Тарихи тұрғыда бұл өнім дәстүрлі мал шаруашылығымен
айналысатын этникалық топтардың күнделікті рационының ажырамас бөлігіне
айналған. Түйе сүті құрамында жоғары сапалы ақуыздармен қатар,
иммуноглобулиндер, лактоферрин, темір, мырыш, С дәрумені сияқты
биологиялық белсенді заттардың болуына байланысты, оны тек тағамдық
емес, сонымен қатар профилактикалық және терапиялық мақсатта қолдану
әлеуеті артып келеді {[}9,10{]}. Ф. Бенмезиан-Дерраджи жүргізген
зерттеулерге сәйкес, сүт өнімдерінің сапасын арттыруда түйе сүтіне
ерекше назар аударады, себебі оның құрамында канцерогендерге қарсы және
иммуностимуляторлық қасиеттерге ие биобелсенді компоненттер --
антиоксиданттар мен лактоферриннің жоғары мөлшері анықталған. Бұл заттар
адам ағзасын патогенді бактериялар мен вирустардың әсерінен қорғауда
маңызды рөл атқарады, сондай-ақ жалпы иммундық жүйенің қызметін
күшейтуге ықпал етеді. {[}11,12{]}.

Сиыр, ешкі және түйе шаруашылығынан алынатын өнімдерді тұтыну --
еліміздің аграрлық секторындағы маңызды құрылымдық элементтердің бірі
болып табылады. Бұл жануарлардан өндірілетін сүт пен сүт өнімдері
халықтың дәстүрлі азық-түлік тұтыну рационында маңызды орынға ие болып,
агроөнеркәсіптік кешеннің тұрақты және тиімді дамуына оң ықпал етеді
{[}13{]}.

{\bfseries Материалдар мен әдістер.} Зерттеу жұмысы Жәңгір хан атындағы
Батыс Қазақстан аграрлық-техникалық университетінің Сынау орталығы
зертханасында орындалды.

Зерттеу нысаны ретінде Қазақстан Республикасының қолданыстағы
стандарттарына (ҚР СТ 1760-2008, ҚР СТ 32259-2013, ҚР СТ 166-2015)
сәйкес келетін сиыр, ешкі және түйе сүттері пайдаланылды. Ашыту процесі
үшін Micromilk өндірушісінің KF көпкомпонентті ұйытқы дақылы қолданылды.
Аталған ұйытқының құрамына \emph{Streptococcus thermophilus, Lactococcus
lactis, Lactococcus cremoris, Leuconostoc mesenteroides subsp. cremoris,
Lactococcus lactis biovar diacetylactis} сияқты сүт қышқылды бактериялар
мен арнайы ашытқылар кіреді.
\end{multicols}

\begin{figs}[1 - сурет. Зерттеу нысандары]
\fig[0.4\textwidth][5cm]{p/image11}
\fig[0.5\textwidth][5cm]{p/image12}
\end{figs}

\begin{multicols}{2}
Айран өндіру технологиясы бойынша сүт үлгілеріне 10:1 қатынасында ұйытқы
қосылып, олар келесідей топтастырылды:

№1 үлгі - сиыр сүті негізіндегі айран;

№2 үлгі - ешкі сүті негізіндегі айран;

№3 үлгі - түйе сүті негізіндегі айран.

Үлгілерді ашыту процесі 33±2 °С температурада, өнімнің титрленетін
қышқылдығы 90-120 °Т көрсеткішіне жеткенге дейін (шамамен 8 сағат
аралығында) жүргізілді.

Зерттеу барысында стандартталған келесі талдау әдістері қолданылды:

Органолептикалық бағалау: ГОСТ ИСО 22935-2-2011 бойынша сыртқы түрі,
консистенциясы, дәмі мен иісі анықталды.

Физика-химиялық көрсеткіштер: майдың массалық үлесі (ГОСТ 5867-90),
ақуыз мөлшері (ГОСТ 25179-90) және лактоза үлесі (ГОСТ 34304-2017)
мемлекетаралық стандарттар негізінде есептелді.

Қышқылдық динамикасы: Титрленетін қышқылдық титриметриялық әдіспен (ГОСТ
3624-92), ал белсенді қышқылдық (рН) рН-150МИ қондырғысында
потенциометриялық әдіспен анықталды.

{\bfseries Нәтижелер мен талқылау.} Зерттеу барысында сиыр, ешкі және түйе
сүтінен KF ұйытқысымен алынған айран үлгілерінің сапалық сипаттамаларын
жан-жақты бағалау мақсатында олардың органолептикалық және
физика-химиялық көрсеткіштеріне талдау жасалды. Органолептикалық бағалау
кезінде өнімдердің сыртқы түрі, иісі, дәмі мен консистенциясы сараланды.
Физика-химиялық параметрлер ретінде майдың, ақуыздың массалық үлестері,
құрғақ заттар мөлшері мен лактоза деңгейі анықталды (кесте 1).

1-кестеде көрсетілгендей, сүттің әр түрінен алынған үлгілердің құрамында
айтарлықтай айырмашылықтар байқалды. Бұл айырмашылықтар тек жануарлардың
түріне ғана емес, сонымен қатар олардың азықтандыру рационына, лактация
кезеңіне және шикізатты технологиялық өңдеу ерекшеліктеріне тікелей
байланысты.

Зерттеу барысында анықталғандай, сиыр және ешкі сүтінен жасалған
айрандар стандартты органолептикалық талаптарға толық сәйкес келсе, түйе
сүтінен алынған үлгінің консистенциясы өзіндік биологиялық құрылымына
(ақуыз және май фракцияларының ерекшелігі) байланысты сұйықтау әрі өзіне
тән ерекше дәмге ие болды. Бұл нәтижелер түйе сүтін ашыту кезінде
технологиялық режимдерді (уақыт пен температураны) жекелей бейімдеу
қажеттігін негіздейді.

Органолептикалық бағалау нәтижелері үш үлгідегі сүттің түсі, иісі, дәмі
және консистенциясы белгіленген стандарттарға сәйкес келетінін көрсетті.
Үлгілер біртектілігімен және жағымды дәмдік қасиеттерімен ерекшеленді,
бұл олардың тұтынушылық тартымдылығын растайды {[}14,15{]}.

Ақуыз бен майдың массалық үлесі өнімнің энергетикалық құндылығы мен
құрылым түзілуіне тікелей әсер етеді. Зерттеу нәтижесінде №2 үлгіде
(ешкі сүті) лактоза мөлшерінің жоғары болуы (5,26\%) анықталды, бұл
ашыту процесінің қарқынды жүруіне қолайлы жағдай жасайды. Түйе сүтінде
(№3 үлгі) ақуыз (3,75\%) мен май (4,88\%) мөлшерінің жоғары болуы оның
тағамдық құндылығын арттырады. Алайда, түйе сүтіндегі ақуыз және құрғақ
заттар (10,05\%) мөлшерінің стандартты талаптардан төмен болуы өнімнің
консистенциясына өзіндік әсерін тигізді.

Түйе сүтінің ашыту процесінің 8 сағатқа дейін созылуы оның биохимиялық
ерекшеліктерімен түсіндіріледі. Түйе сүтіндегі казеин фракцияларының
құрылымы (әсіресе β-казеиннің жоғары үлесі және κ-казеиннің аздығы) мен
мицеллалардың үлкен мөлшері тығыз ұйытындының (ұйытқының) кеш түзілуіне
себеп болады. Сондай-ақ, түйе сүтінің құрамындағы лактоферрин мен
иммуноглобулиндер сияқты табиғи антимикробтық заттардың жоғары
концентрациясы ашыту бактерияларының белсенділігін тежеп, қышқылдықтың
баяу көтерілуіне ықпал етеді. Бұл факторлар түйе сүтінен айран алу
кезінде технологиялық процесті (ашыту уақытын ұзарту немесе
температураны оңтайландыру) түзетуді талап етеді.

Ашыту процесі 8 сағат бойы бақыланды. Қышқыл түзілу динамикасын зерттеу
мақсатында сынамалар әр сағат сайын алынып отырды. Барлық үлгілерде рН
мәнінің төмендеуі және титрленетін қышқылдықтың артуы байқалды, бұл KF
ұйытқысының сүт ортасына бейімделіп, лактозаны сүт қышқылына айналдыру
процесінің жүріп жатқанын дәлелдейді {[}16{]}.
\end{multicols}

\tcap{1 - кесте. Сиыр, ешкі және түйе сүттерінің органолептикалық және физика-химиялық көрсеткіштері}
\begin{longtblr}[
  label = none,
  entry = none,
  note{} = {\emph{*көрсеткіштерге қатысты ҚР СТ және ГОСТ талаптарыгда белгіленбеген}},
]{
  width = \linewidth,
  colspec = {Q[100]Q[360]Q[250]Q[152]},
  cells = {c},
  cells = {font = \small},
  row{1} = {font=\bfseries\small},
  cell{1}{1} = {c=2}{},
  cell{2}{1} = {r=7}{},
  cell{9}{1} = {r=7}{},
  cell{16}{1} = {r=7}{},
  hlines,
  vlines,
}
Көрсеткіштер                  &                                       & Сипаттамасы және нормасы                                        & Зерттеу нәтижелері \\
{Сиыр сүті\\ҚР СТ\\1760-2008} & Дәмі мен иісі                         & Таза, балғын табиғи сүтке тән емес бөгде иістер мен дәмдерсіз   & Сәйкес             \\
                              & Консиситенциясы                       & Біртекті, шөгіндісіз                                            & Сәйкес             \\
                              & Түсі                                  & Ақтан сәл сарыға дейін                                          & Сәйкес             \\
                              & Майдың массалық үлесі, \%, кем емес   & 3,2                                                             & 4,06±0,001         \\
                              & Ақуыздың массалық үлесі, \%, кем емес & 2,8                                                             & 3,23±0,003         \\
                              & Құрғақ заттар, \%*                    & -                                                               & 8,95±0,001         \\
                              & Лактоза, \%*                          & -                                                               & 3,41±0,002         \\
{Ешкі сүті\\ГОСТ 32259-2013}  & Дәмі мен иісі                         & Ерекше ешкі сүтінің дәмі, ешкі сүтіне тән иіс                   & Сәйкес             \\
                              & Консиситенциясы                       & Сұйық әрі біркелкі                                              & Сәйкес             \\
                              & Түсі                                  & Ақ немесе ақшыл-сары                                            & Сәйкес             \\
                              & Майдың массалық үлесі, \%             & 2,8-5,6                                                         & 5,01±0,002         \\
                              & Ақуыздың массалық үлесі, \%, кем емес & 3,0                                                             & 3,50±0,001         \\
                              & Құрғақ заттар, \%*                    & -                                                               & 9,59±0,001         \\
                              & Лактоза, \%*                          & -                                                               & 5,26±0,002         \\
{Түйе сүті\\ҚР СТ\\166-2015}  & Дәмі мен иісі                         & Таза, жаңа сауылган сүтке тән емес бөгде дәмдер мен иістері жоқ & Сәйкес             \\
                              & Консиситенциясы                       & Біртекті, шөгіндісіз және қауысыз                               & Сәйкес             \\
                              & Түсі                                  & Ақтан ақшыл сарығы дейін                                        & Сәйкес             \\
                              & Майдың массалық үлесі, \%, кем емес   & 3,0                                                             & 4,88±0,001         \\
                              & Ақуыздың массалық үлесі, \%, кем емес & 3,8                                                             & 3,75±0,003         \\
                              & Құрғақ заттектердің орта мелшері, \%  & 15,0                                                            & 10,05±0,003        \\
                              & Лактоза, \%*                          & -                                                               & 4,45±0,001         
\end{longtblr}

\begin{multicols}{2}
2-суретте көрсетілген график бойынша 8 сағат ашыту кезінде титрленетін
қышқылдық 90-120 °T аралығындағы мәндерге жететінін көруге болады.
Титрленетін қышқылдықтың ашыту ұзақтығына тәуелділігі бойынша ең
белсенді сүттерге №1 және №2 үлгілер, яғни сиыр мен ешкінің сүтін
жатқызуға болады. Бұл үлгілерде қышқылдықтың қарқынды артуы KF
ұйытқысының құрамындағы термофилді стрептококтардың жоғары ферментативті
белсенділігін көрсетеді. Ал № 3 үлгі - түйе сүтінде қышқылдық мәнінің
баяу артатыны белгілі болды, бұл оның жоғары буферлік сыйымдылығымен
және антимикробтық белоктардың әсерімен байланысты.

Сүттің белсенді қышқылдығы (рН) бұл ерітіндідегі сутегі иондарының
концентрациясын көрсететін маңызды көрсеткіш. Ол сүттегі қышқылдардың
мөлшерін білдіреді және сүттің рН мәнін анықтайды. Белсенді қышқылдық,
көбінесе сүттің қышқылдануы немесе ашуы кезінде байқалады. Мұндай
өзгерістер көбінесе сүтте сүт қышқылы бактерияларының әрекет етуінен,
атап айтқанда лактозаның сүт қышқылына айналуынан пайда болады. Белсенді
қышқылдықты өлшеу арқылы сүттің сапасын бағалауға болады, өйткені рН
мәнінің күрт төмендеуі микробиологиялық процестердің белсенділігін және
өнімнің ұю (коагуляция) нүктесіне жеткенін білдіреді {[}17{]}.

Зерттеу барысында түйе сүтінің рН көрсеткіші сиыр және ешкі сүтіне
қарағанда баяу төмендегені байқалды. Бұл технологиялық тұрғыдан түйе
сүтін ашыту кезінде термостаттау уақытын ұзартуды (8 сағатқа дейін)
немесе ұйытқы мөлшерін арттыру қажеттілігін негіздейді.
\end{multicols}

\fig[0.6\textwidth]{p/image43}[2 - сурет. Ашыту динамикасы]
\fig[0.55\textwidth]{p/image44}[3 - сурет. Белсенді қышқылдық динамикасы (рН)]

\begin{multicols}{2}
3-суретте берілген белсенді қышқылдықтың ашыту ұзақтығына тәуелділік
графигі бойынша KF ұйытқысының зерттелген сүттің барлық үш түрінде де
белсенді жұмыс жасайтыны дәлелденді. Атап айтқанда, № 2 үлгіде (ешкі
сүті) белсенді қышқылдықтың (рН) артуы неғұрлым қарқынды жүрді, бұл ешкі
сүтінің микрофлораның дамуы үшін қолайлы биотехнологиялық орта екенін
көрсетеді. Ал түйе сүтінде (№ 3 үлгі) белсенді қышқылдық мәнінің өзгеруі
баяу жүреді.

Зерттеу нәтижелеріне сүйене отырып, үш сағаттан астам ұзақ өсіру кезінде
титрленетін қышқылдық сүт қышқылды микроорганизмдердің өсуіне қолайлы
температуралық жағдай (33±2 °С) жасау арқылы тұрақты артады деген
қорытынды жасауға болады. Қышқылдықтың күрт өсуі 4-5 сағат аралығында
байқалады, бұл кезең микрофлораның логарифмдік өсу фазасы үшін оңтайлы
болып табылады.

Түйе сүтіндегі ашыту процесінің 8 сағатқа дейін созылуы мен қышқылдық
деңгейінің баяу көтерілуі оның ерекше биохимиялық құрамымен ---
құрамындағы антимикробтық белоктардың (лактоферрин, лизоцим) жоғары
мөлшерімен және казеин мицеллаларының құрылымымен байланысты. Бұл
факторлар тромбтың (ұйытындының) кеш түзілуіне әсер еткенімен, KF
ұйытқысы 8 сағат ішінде өнімнің қажетті технологиялық кондицияға жетуін
қамтамасыз етті.

Зерттеу барысында сиыр, ешкі және түйе сүтінен дайындалған айранның
органолептикалық көрсеткіштері (сыртқы түрі, иісі, дәмі,
консистенциясы), сондай-ақ физика-химиялық көрсеткіштері (май, ақуыз,
құрғақ заттар және лактоза) анықталып, 2-кестеде жүйеленді (кесте 2).
\end{multicols}

\tcap{2 - кесте. Micromilk ұйытқысымен дайындалған айрандардың органолептикалық көрсеткіштері}
\begin{longtblr}[
  label = none,
  entry = none,
  note{} = {\emph{*көрсеткіштерге қатысты ГОСТ талаптарыгда белгіленбеген}},
]{
  width = \linewidth,
  colspec = {Q[160]Q[400]Q[100]Q[100]Q[100]},
  cells = {c},
  cells = {font = \small},
  hlines,
  vlines,
}
Көрсеткіштер                          & ГОСТ 31702-2013. Айран. Техникалық шарттары                                                                                                  & Үлгі №1     & Үлгі №2     & Үлг №3      \\
Дәмі мен иісі                         & Таза, сүтқышқылды, бөгде дәмдер мен иістерсіз немесе тұз бен су қосылған кезде тұзды                                                         & Сәйкес      & Сәйкес      & Сәйкес      \\
Консистенциясы                        & Біртекті, Араластырған нан кейін жоғалып кететін сарысудың болуына рұқсат етіледі. Қалыпты микрофлорадан туындаған жеке көздер түріндегі газ & Сәйкес      & Сәйкес      & Сәйкес      \\
Түсі                                  & Сүтті ақ, бүкіл массасы бойынша біркелкі                                                                                                     & Сәйкес      & Сәйкес      & Сәйкес      \\
Майдың массалық үлесі, \%, кем емес   & 0,5-4,0                                                                                                                                      & 3,88±0,002  & 4,15±0,001  & 3,93±0,003  \\
Ақуыздың массалық үлесі, \%, кем емес & 2,8                                                                                                                                          & 4,12±0,001  & 4,31±0,01   & 3,94±0,001  \\
Құрғақ заттар, \%*                    & - ⁕                                                                                                                                          & 12,90±0,002 & 13,15±0,001 & 12,58±0,003 \\
Лактоза, \%*                          & - ⁕                                                                                                                                          & 3,35±0,002  & 3,31±0,003  & 3,48±0,002  
\end{longtblr}

\begin{multicols}{2}
Кесте 2-ге сәйкес, KF ұйытқысымен дайындалған айранның органолептикалық
сипаттамалары (дәмі, хош иісі, консистенциясы және түсі) пайдаланылған
шикізат түріне тікелей байланысты. Сиыр сүтінен (№ 1 үлгі) алынған айран
классикалық, жұмсақ ашытылған сүт дәмімен және біртекті
консистенциясымен ерекшеленді. Ешкі сүтінен (№ 2 үлгі) алынған өнімде
айқын әрі ерекше спецификалық дәм байқалды, бұл ешкі сүтінің май
қышқылдық құрамымен және KF микрофлорасының метаболизмімен байланысты.

Түйе сүтінен жасалған айран (№ 3 үлгі) сұйық консистенциясымен және
жеңіл қышқыл дәмімен ерекшеленеді. Зерттеу нәтижелері көрсеткендей, түйе
сүтінде ақуыз және құрғақ заттар мөлшерінің төмен болуы ұйытындының
тығыздығына тікелей әсер еткен. Дегенмен, KF ұйытқысының құрамындағы
ашытқылар түйе сүтіне тән дәмді жұмсартып, өнімге жағымды хош иіс берді.

Өнімдердің сенсорлық сипаттамаларын жүйелі түрде салыстыру мақсатында
органолептикалық профилограмма құрастырылды (сурет 4). Профилограмма
дәмін, иісін, түсін және консистенциясын сандық (баллдық) әдіспен
бағалауға мүмкіндік берді {[}18{]}.

Профилограмма нәтижесі көрсеткендей, № 2 үлгі (ешкі сүті) дәм мен хош
иіс айқындылығы бойынша ең жоғары балл алса, № 1 үлгі (сиыр сүті)
консистенцияның тұрақтылығы мен біртектілігі жағынан көш бастады. Түйе
сүтінен алынған айран консистенция бойынша төменірек балл алғанымен,
өзінің диеталық және емдік қасиеттеріне тән ерекше дәмдік профилімен
жоғары бағаланды. Бұл субъективті бағалауларды стандартталған формада
ұсыну өнімнің нарықтағы тартымдылығын ғылыми негіздеуге мүмкіндік
береді.
\end{multicols}

\fig[0.5\textwidth]{p/image45}[4 - сурет. Зерттеу үлгілерінің дегустациялық сараптамасы]

\begin{multicols}{2}
4-суретте көрсетілген органолептикалық профилограмма (5 баллдық жүйе
бойынша) KF ұйытқысымен ашытылған айран үлгілерінің сенсорлық
сипаттамаларын салыстыруға мүмкіндік береді. Сараптама нәтижелері
көрсеткендей, сиыр және ешкі сүтінен дайындалған айрандар (№ 1 және № 2
үлгілер) дәмдік қасиеттері мен консистенциясы бойынша ең жоғары балл
жинады.

Дәмі мен хош иісі: Ешкі сүтінен алынған айран өзінің айқын әрі жағымды
спецификалық дәмімен ерекшеленсе, сиыр сүтінен жасалған үлгі классикалық
ашытылған сүт өніміне тән жұмсақ дәм көрсетті. Түйе сүтіндегі (№3 үлгі)
түйе сүтіне тән өзіндік иіс пен дәм сақталып қалғаны байқалды, бұл
шикізаттың табиғи қасиетінің жоғары болуымен түсіндіріледі.

Консистенция: № 1 және № 2 үлгілер біртекті, тығыз және тұтқыр
консистенциясы үшін жоғары бағаланды. Бұл KF ұйытқысының сиыр және ешкі
сүтінің ақуыздарымен жақсы әрекеттесіп, тұрақты құрылым түзетінін
дәлелдейді. Түйе сүтінен алынған айранның консистенциясы сұйықтау болды,
бұл оның құрамындағы құрғақ заттардың төмендігімен және казеин
фракцияларының ерекшелігімен байланысты.

Осылайша, KF ұйытқысы сиыр және ешкі сүтінен жоғары органолептикалық
сапасы бар өнімдер алуға мүмкіндік береді. Түйе сүтінен алынған айранның
сұйық консистенциясына қарамастан, оның өзіндік дәмі мен иісі
функционалдық мақсаттағы өнім ретінде жоғары бағалануына негіз болады.

{\bfseries Қорытынды.} Зерттеу нәтижелері айранның органолептикалық және
физика-химиялық сипаттамаларына пайдаланылатын сүт түрінің елеулі ықпал
ететінін көрсетті.

KF ұйытқысының тиімділігі: Micromilk (KF) ұйытқысының зерттелген сүт
түрлерінің (сиыр, ешкі, түйе) барлығында сүт қышқылды бактериялар мен
ашытқылардың белсенді дамуын қамтамасыз ететіні расталды. Дайын
өнімдердің титрленетін қышқылдығы 90-120 °Т шегінде болып, сапалы
ашытылған өнім стандартына сәйкес келді.

Шикізат ерекшеліктері: Сиыр сүтінен дайындалған айран классикалық
консистенциясымен ерекшеленсе, ешкі сүтінен алынған үлгі ақуыз (3,50\%)
бен майдың (5,01\%) жоғары құрамымен және ерекше дәмдік профилімен
бағаланды.

Түйе сүтінің технологиялық ерекшелігі: Түйе сүтіндегі ашыту процесінің 8
сағатқа дейін созылуы оның биохимиялық құрамымен --- құрғақ заттардың
төмендігімен (10,05\%) және казеин фракцияларының ерекшеліктерімен
негізделді. Түйе сүтіндегі табиғи антимикробтық белоктардың
(лактоферрин, иммуноглобулин) тежегіш әсері тромбтың (ұйытындының)
кешігуіне себеп болады. Осыған байланысты түйе сүті үшін оңтайлы ашыту
режимі 33±2 °С температурада 8 сағат деп белгіленді.

Практикалық маңызы: Алынған нәтижелер әртүрлі шикізат түрлерінің
ерекшеліктерін ескере отырып, айран өндіру технологиясын оңтайландыруға
және жақсартылған функционалдық қасиеттері бар жаңа ашытылған сүт
өнімдерін әзірлеуге ғылыми-тәжірибелік негіз болады.
\end{multicols}

\begin{center}
{\bfseries Әдебиеттер}
\end{center}

\begin{refs}
1. Nguyen A., Francis M., Windfeld E., Lhermie G., Kim K. Developing an
immersive virtual farm simulation for engaging and effective public
education about the dairy industry // Computers \& Graphics.-
2024. -Vol..118.-P.173-183. .

2. Канарейкина С.Г. Комбинированный продукт с использованием сухого
кобыльего молока // Коневодство и конный спорт. -2014. -№.2. - С.
29-31.

3. Асенова Б.К. Показатели качества и идентификации кефира //
Международный научно-исследовательский журнал. - 2015. - № 6(37). - C.
13-17.

4. Оспанов А.Б. Оценка возможности использования козьего и овечьего
молока в производстве йогуртов // Ползуновский вестник. - 2022. -Т.4(1)-
С.154-159.

DOI 10.25712/ASTU.2072-8921.2022.04.020.

5. Paszczyk B., Czarnowska-Kujawska M. Paszczyk B., Klepacka J., Tońska
E. Health-promoting ingredients in goat's milk and fermented goat's milk
drinks //Animals. - 2023.-Vol.13 (5):907.
\href{https://doi.org/10.3390/ani13050907}{DOI 10.3390/ani13050907}.

6. Садовой В.В., Вобликова Т.В., Пермяков А.В. Жирнокислотный состав
козьего и овечьего молока и его трансформация в процессе производства
йогурта // Техника и технология пищевых производств.- 2019.-Т.49(4).- С.
555-562. DOI 10.21603/2074-9414-2019-4-555-562.

7. Шувариков А.С. Канина К.А., Красуля О.Н., Пастух О.Н., Робкова Т.О.
Физико-химические показатели козьего, овечьего и коровьего молока //
Овцы, козы, шерстяное дело.-2017.-№ 1. -С.38-40.

8. Singhal S. Baker R.D., Baker S.S. A comparison of the nutritional
value of cow' s milk and nondairy beverages // Journal of
pediatric gastroenterology and nutrition.-2017.-Vol.64 (5). - P.
799-805. ~DOI
~\href{https://doi.org/10.1097/mpg.0000000000001380}{10.1097/MPG.0000000000001380}.

9. Суюнчев О.А., Вобликова Т.В. Сыры из козьего молока // Сыроделие и
маслоделие. - 2007.-№ 5.- С.10-11.

10. Konuspayeva G., Faye B. Recent advances in camel milk processing
//Animals.-2021.-Vol.11 (4):1045.
\href{https://doi.org/10.3390/ani11041045}{DOI 10.3390/ani11041045}.

11. Лысенко К. Пикантный йогурт // Переработка молока. -2017.-№ 5 (212).
-С.44-45.

12. Benmeziane--Derradji F. Evaluation of camel milk: gross compositiona
scientific overview // Tropical Animal Health and
Production.-2021.-Vol.53 (2): 308. DOI
~\href{https://doi.org/10.1007/s11250-021-02689-0}{10.1007/s11250-021-02689-0}.

13. Getaneh G., Mebrat A., Wubie A., Kendie H. Review on goat milk
composition and its nutritive value // Journal of Nutrition and Health
Sciences.-2016.-Vol.3 (4).-P.1-10. DOI
\href{http://dx.doi.org/10.15744/2393-9060.3.401}{10.15744/2393-9060.3.401}.

14. Lad S.S., Aparnathi K.D., Mehta B., Velpula S. Goat milk in human
nutrition and health--a review // Int. J. Curr. Microbiol. Appl.
Sci.-2017.-Vol.6(5).-P.1781-1792. DOI
\href{http://dx.doi.org/10.20546/ijcmas.2017.605.194}{10.20546/ijcmas.2017.605.194}.

15. Yeleupaeva Sh.K. Zhumagalyeva Zh.Zh., Zhuzbaeva G.O. Bio chemical
properties of sour-milk bacteria in preparation of drinks //
International Scientific and Practical Conference World
Science.-2017.-Vol.5 (4).- P.44-45.

16. Seifu E. Camel milk products: innovations, limitations and
opportunities //Food Production, Processing and Nutrition. - 2023. -
Vol.5 (1). - P.15. DOI 10.1186/s43014-023-00130-7.

17. Ho T.M., Zou Z., Bansal N. Camel milk: A review of its nutritional
value, heat stability, and potential food products//Food Research
International.-2022.-Vol.153:110870.
\href{https://doi.org/10.1016/j.foodres.2021.110870}{DOI
10.1016/j.foodres.2021.110870}.

18. Dan T., Hu H., Li T., Dai A., He B., Wang Y. Screening of
mixed-species starter cultures for increasing flavour during
fermentation of milk//International Dairy Journal.-2022.-Vol.135:
105473. \href{https://doi.org/10.1016/j.idairyj.2022.105473}{DOI
10.1016/j.idairyj.2022.105473}.
\end{refs}

\begin{center}
{\bfseries References}
\end{center}

\begin{refs}
1. Nguyen A., Francis M., Windfeld E., Lhermie G., Kim K. // Developing
an immersive virtual farm simulation for engaging and effective public
education about the dairy industry // Computers \& Graphics.-
2024. -Vol..118.-P.173-183. .

2. Kanarejkina S.G. Kombinirovannyj produkt s
ispol' zovaniem suhogo kobyl' ego moloka
// Konevodstvo i konnyj sport. -2014.-№.2. - S.29-31. {[}in Russian{]}

3. Asenova B.K. Pokazateli kachestva i identifikacii kefira //
Mezhdunarodnyj nauchno-issledovatel' skij zhurnal.-
2015. - № 6(37). - C.13-17. {[}in Russian{]}

4. Ospanov A.B. Ocenka vozmozhnosti ispol' zovanija
koz' ego i ovech' ego moloka v
proizvodstve jogurtov // Polzunovskij vestnik. - 2022. -T.4(1)-
S.154-159.

DOI 10.25712/ASTU.2072-8921.2022.04.020. {[}in Russian{]}

5. Paszczyk B., Czarnowska-Kujawska M. Paszczyk B., Klepacka J., Tońska
E. Health-promoting ingredients in goat's milk and fermented goat's milk
drinks //Animals. - 2023.-Vol.13 (5):907.
\href{https://doi.org/10.3390/ani13050907}{DOI 10.3390/ani13050907}.

6. Sadovoj V.V., Voblikova T.V., Permjakov A.V. Zhirnokislotnyj sostav
koz' ego i ovech' ego moloka i ego
transformacija v processe proizvodstva jogurta // Tehnika i tehnologija
pishhevyh proizvodstv.- 2019.-T.49(4).- S.555-562. DOI
10.21603/2074-9414-2019-4-555-562. {[}in Russian{]}

7. Shuvarikov A.S. Kanina K.A., Krasulja O.N., Pastuh O.N., Robkova T.O.
Fiziko-himicheskie pokazateli koz' ego,
ovech' ego i korov' ego moloka // Ovcy,
kozy, sherstjanoe delo.-2017.-№ 1. -S.38-40. {[}in Russian{]}

8. Singhal S. Baker R.D., Baker S.S. A comparison of the nutritional
value of cow' s milk and nondairy beverages // Journal of
pediatric gastroenterology and nutrition.-2017.-Vol.64 (5). - P.
799-805. ~DOI
~\href{https://doi.org/10.1097/mpg.0000000000001380}{10.1097/MPG.0000000000001380}.

9. Sujunchev O.A., Voblikova T.V. Syry iz koz' ego moloka
// Syrodelie i maslodelie. - 2007.-№ 5.- S.10-11. {[}in Russian{]}

10. Konuspayeva G., Faye B. Recent advances in camel milk processing
//Animals.-2021.-Vol.11 (4):1045.
\href{https://doi.org/10.3390/ani11041045}{DOI 10.3390/ani11041045}.

11. Lysenko K. Pikantnyj jogurt // Pererabotka moloka. -2017.-№ 5
(212).-S.44-45. {[}in Russian{]}

12. Benmeziane--Derradji F. Evaluation of camel milk: gross composition:
a scientific overview // Tropical Animal Health and
Production.-2021.-Vol.53 (2): 308. DOI
~\href{https://doi.org/10.1007/s11250-021-02689-0}{10.1007/s11250-021-02689-0}.

13. Getaneh G., Mebrat A., Wubie A., Kendie H. Review on goat milk
composition and its nutritive value//Journal of Nutrition and Health
Sciences.-2016.-Vol.3 (4).-P.1-10. DOI
\href{http://dx.doi.org/10.15744/2393-9060.3.401}{10.15744/2393-9060.3.401}.

14. Lad S.S., Aparnathi K.D., Mehta B., Velpula S. Goat milk in human
nutrition and health--a review // Int. J. Curr. Microbiol. Appl.
Sci.-2017.-Vol.6(5).-P.1781-1792. DOI
\href{http://dx.doi.org/10.20546/ijcmas.2017.605.194}{10.20546/ijcmas.2017.605.194}.

15. Yeleupaeva Sh.K. Zhumagalyeva Zh.Zh., Zhuzbaeva G.O. Biochemical
properties of sour-milk bacteria in preparation of drinks //
International Scien tific and Practical Conference World
Science.-2017.-Vol.5 (4).- P.44-45.

16. Seifu E. Camel milk products: innovations, limitations and
opportunities // Food Production, Processing and Nutrition.- 2023.-
Vol.5 (1). - P.15. DOI 10.1186/s43014-023-00130-7.

17. Ho T.M., Zou Z., Bansal N. Camel milk: A review of its nutritional
value, heat stability, and potential food products// Food Research
International.-2022.-Vol.153:110870.
\href{https://doi.org/10.1016/j.foodres.2021.110870}{DOI
10.1016/j.foodres.2021.110870}.

18. Dan T., Hu H., Li T., Dai A., He B., Wang Y. Screening of
mixed-species starter cultures for increasing flavour during
fermentation of milk// International Dairy Journal.-2022.-Vol.135:
105473. \href{https://doi.org/10.1016/j.idairyj.2022.105473}{DOI
10.1016/j.idairyj.2022.105473}.
\end{refs}

\begin{info}
\hspace{1em}\emph{{\bfseries Авторлар туралы мәлімет}}

Оразов А.Ж. - т.ғ.к., доцент м.а. Жәңгір хан атындағы Батыс Қазақстан
аграрлық-техникалық университеті, Орал, Қазақстан, е-mail:
orazov\_ayan@mail.ru;

Рыскалиева Б.Ж..- магистр, аға оқытушы, Жәңгір хан атындағы Батыс
Қазақстан аграрлық-техникалық университеті, Орал, Қазақстан, е-mail:
bryskalieva@mail.ru;

Сармалаев А.А.-магистр, оқытушы, Жәңгір хан атындағы Батыс Қазақстан
аграрлық-техникалық университеті, Орал, Қазақстан, е-mail:
sarmalayev@gmail.com;

Меркенова А.М. - магистрант, Жәңгір хан атындағы Батыс Қазақстан
аграрлық-техникалық университеті, Орал, Қазақстан, е-mail:
merkenova@bk.ru.

\hspace{1em}\emph{{\bfseries Information about the authors}}

Orazov A. Zh. - candidate of Technical Sciences, acting associate
professor Zhangir Khan West Kazakhstan agrarian and technical
University, Uralsk, Kazakhstan, e-mail:
orazov\_ayan@mail.ru;

Ryskaliyeva B. Zh., - master' s degree, senior lecturer,
West Kazakhstan agrarian and technical University named after Zhangir
Khan, Uralsk, Kazakhstan, e-mail:
bryskalieva@mail.ru;

Sarmalaev A.A., - master' s degree, Teacher, West
Kazakhstan agrarian and technical University named after Zhangir Khan,
Uralsk, Kazakhstan, e-mail:
sarmalayev@gmail.com;

Merkenova A.M. - master' s student, West Kazakhstan
agrarian and technical University named after Zhangir Khan, Uralsk,
Kazakhstan, e-mail:
merkenova@bk.ru.
\end{info}
